\usepackage{listings}
\usepackage{listingsutf8}
\usepackage{pifont}
\usepackage{tcolorbox}
\usepackage{setspace}


\renewcommand{\baselinestretch}{1.1} 

%\usepackage{animate}
\newcommand{\insertvideo}[3]{\begin{flushright}\fbox{Vidéo #3: \sphinxurl{#2}  \includegraphics[scale=0.3]{#1}}\end{flushright}}
\newcommand{\insertsolution}[1]{\flushright \rotatebox[origin=c]{180}{\textit{Réponse: #1}}}
\newcommand{\glossaryterm}[1]{\textit{#1}}
\newcommand{\glossarylist}{\\}
\renewenvironment{wrapfigure}[2]{\begin{figure}[htbp]}{\end{figure} \ignorespacesafterend}
\newenvironment{timeline}{}{}
\newenvironment{timelineitem}[2]{}{}

% reduce outside margin
\newlength{\margindiff}
\setlength{\margindiff}{1em}
\addtolength{\oddsidemargin}{\margindiff}
\addtolength{\evensidemargin}{-\margindiff}
\addtolength{\textwidth}{\margindiff}

%% default admonition type
%\renewenvironment{sphinxnote}[1]{\begin{tcolorbox}[colframe=blue]\sphinxstrong{#1}}{\end{tcolorbox}}

\definecolor{babyblue}{rgb}{0.54, 0.81, 0.94}
\definecolor{llgray}{rgb}{0.96, 0.96, 0.96}
\definecolor{carolinablue}{rgb}{0.6, 0.73, 0.89}
\definecolor{blizzardblue}{rgb}{0.67, 0.9, 0.93}
\definecolor{beaublue}{rgb}{0.74, 0.83, 0.9}
\definecolor{brightturquoise}{rgb}{0.03, 0.91, 0.87}
\definecolor{darkviolet}{rgb}{0.58, 0.0, 0.83}
\definecolor{darkspringgreen}{rgb}{0.09, 0.45, 0.27}
\definecolor{lapislazuli}{rgb}{0.15, 0.38, 0.61}
\definecolor{ultraviolet}{RGB}{89, 51, 209}
\definecolor{ultraturquoise}{RGB}{92, 201, 192}
\definecolor{bondiblue}{rgb}{0.0, 0.58, 0.71}
\definecolor{neoncarrot}{rgb}{1.0, 0.64, 0.26}
\definecolor{radicalred}{rgb}{1.0, 0.21, 0.37}
\definecolor{orangemodulo}{RGB}{241, 180, 63}


%% specific admonition type - use tcolorbox arguments (or something else) to adapt latex rendering
% TODO: this is a problem with page breaks, we may want to use activitybox instead
% (cf. https://tex.stackexchange.com/questions/350316/create-manual-page-break-within-tcolorbox)
\newcommand\spacebeforeadmonition{\vspace{10pt}}
\newcommand\spaceafteradmonition{\vspace{3pt}}
\newcommand\admonitiontitle[1]{{\fontfamily{phv}\textbf{#1}}}
\newenvironment{exercise}[1]{\spacebeforeadmonition\begin{tcolorbox}[parbox=false, colframe=ultraviolet, colback=llgray, boxrule=1pt, top=8pt, left=8pt, right=28pt, bottom=8pt]\admonitiontitle{#1}}{\end{tcolorbox}\spaceafteradmonition}
\newenvironment{solution}[1]{\spacebeforeadmonition\begin{tcolorbox}[parbox=false, colframe=ultraturquoise, colback=llgray, boxrule=1pt, top=8pt, left=8pt, right=28pt, bottom=8pt]\admonitiontitle{#1}}{\end{tcolorbox}\spaceafteradmonition}
\newenvironment{microactivity}[1]{\spacebeforeadmonition\begin{tcolorbox}[parbox=false, colframe=ultraviolet, colback=llgray, boxrule=1pt, top=8pt, left=8pt, right=28pt, bottom=8pt]\admonitiontitle{#1}}{\end{tcolorbox}\spaceafteradmonition}
\newenvironment{togofurther}[1]{\spacebeforeadmonition\begin{tcolorbox}[parbox=false, colframe=orangemodulo, colback=llgray, boxrule=1pt, top=8pt, left=8pt, right=28pt, bottom=8pt]\admonitiontitle{#1}}{\end{tcolorbox}\spaceafteradmonition}
\newenvironment{important}[1]{\spacebeforeadmonition\begin{tcolorbox}[parbox=false, colframe=radicalred, colback=llgray, boxrule=1pt, top=8pt, left=8pt, right=28pt, bottom=8pt]\admonitiontitle{#1}}{\end{tcolorbox}\spaceafteradmonition}
\newenvironment{didyouknow}[1]{\spacebeforeadmonition\begin{tcolorbox}[parbox=false, colframe=orangemodulo, colback=llgray, boxrule=1pt, top=8pt, left=8pt, right=28pt, bottom=8pt]\admonitiontitle{#1}}{\end{tcolorbox}\spaceafteradmonition}
\newenvironment{reminder}[1]{\spacebeforeadmonition\begin{tcolorbox}[parbox=false, colframe=radicalred, colback=llgray, boxrule=1pt, top=8pt, left=8pt, right=28pt, bottom=8pt]\admonitiontitle{#1}}{\end{tcolorbox}\spaceafteradmonition}
\newenvironment{related}[1]{\spacebeforeadmonition\begin{tcolorbox}[parbox=false, colframe=neoncarrot, colback=llgray, boxrule=1pt, top=8pt, left=8pt, right=28pt, bottom=8pt]\admonitiontitle{#1}}{\end{tcolorbox}\spaceafteradmonition}
\newenvironment{evaluation}[1]{\spacebeforeadmonition\begin{tcolorbox}[parbox=false, colframe=radicalred, colback=llgray, boxrule=1pt, top=8pt, left=8pt, right=28pt, bottom=8pt]\admonitiontitle{#1}}{\end{tcolorbox}\spaceafteradmonition}
\newenvironment{thinkingmatter}[1]{\spacebeforeadmonition\begin{tcolorbox}[parbox=false, colframe=ultraviolet, colback=llgray, boxrule=1pt, top=8pt, left=8pt, right=28pt, bottom=8pt]\admonitiontitle{#1}}{\end{tcolorbox}\spaceafteradmonition}
\newenvironment{note}[1]{\spacebeforeadmonition\begin{tcolorbox}[parbox=false, colframe=neoncarrot, colback=llgray, boxrule=1pt, top=8pt, left=8pt, right=28pt, bottom=8pt]\admonitiontitle{#1}}{\end{tcolorbox}\spaceafteradmonition}
\newenvironment{torecall}[1]{\spacebeforeadmonition\begin{tcolorbox}[parbox=false, colframe=radicalred, colback=llgray, boxrule=1pt, top=8pt, left=8pt, right=28pt, bottom=8pt]\admonitiontitle{#1}}{\end{tcolorbox}\spaceafteradmonition}
\newenvironment{historicdocument}[1]{\spacebeforeadmonition\begin{tcolorbox}[colframe=neoncarrot, colback=llgray, boxrule=1pt, top=8pt, left=8pt, right=28pt, bottom=8pt]\admonitiontitle{#1}}{\end{tcolorbox}\spaceafteradmonition}

% set toc links color to black
\hypersetup{linkcolor=black}


\sphinxpxdimen=.5bp\relax

\definecolor{light-gray}{gray}{0.95} 
\lstset{language=Python,upquote=true,showstringspaces=false,columns=fullflexible, basicstyle = \ttfamily,backgroundcolor = \color{light-gray},numbers=right,stepnumber=1,keywordstyle=\color{blue},stringstyle=\color{magenta},commentstyle=\color{red},breaklines=true}
%% for accented characters in lstlisting
\lstset{inputencoding=utf8/latin1, literate=%
         {à}{{\`a }}1
         {â}{{\^a }}1
         {é}{{\'e}}1
         {è}{{\`e}}1
         {ê}{{\^e}}1
         {ô}{{\^o}}1
         {’}{{'}}1
         {ù}{{\`u}}1
         {û}{{\^u}}1
         {À}{{\`A}}1
         {♣}{{\ding{168}}}1
         {♦}{{\ding{169}}}1
         {♥}{{\ding{170}}}1
         {♠}{{\ding{171}}}1
}

%% undefined unicode characters
\DeclareUnicodeCharacter{270F}{}
\DeclareUnicodeCharacter{FE0F}{}
\DeclareUnicodeCharacter{1F4D2}{}
\DeclareUnicodeCharacter{1F50C}{}
\DeclareUnicodeCharacter{1F1E7}{}
\DeclareUnicodeCharacter{1F1EB}{}
\DeclareUnicodeCharacter{1F1EC}{}
\DeclareUnicodeCharacter{1F1EE}{}
\DeclareUnicodeCharacter{1F1F1}{}
\DeclareUnicodeCharacter{1F1F3}{}
\DeclareUnicodeCharacter{1F1FA}{}
\DeclareUnicodeCharacter{1F1F8}{}
\DeclareUnicodeCharacter{1F1FF}{}
\DeclareUnicodeCharacter{2009}{}
\DeclareUnicodeCharacter{200A}{}
\DeclareUnicodeCharacter{202F}{}
\DeclareUnicodeCharacter{2248}{$\approx$}
\DeclareUnicodeCharacter{300}{è}
\DeclareUnicodeCharacter{301}{é}
\DeclareUnicodeCharacter{302}{ê}
\DeclareUnicodeCharacter{03B1}{$\alpha$}
\DeclareUnicodeCharacter{03B2}{$\beta$}
\DeclareUnicodeCharacter{03B3}{$\gamma$}
\DeclareUnicodeCharacter{03B4}{$\delta$}
\DeclareUnicodeCharacter{03B5}{$\epsilon$}
\DeclareUnicodeCharacter{03B7}{$\eta$}
\DeclareUnicodeCharacter{03B9}{$\iota$}
\DeclareUnicodeCharacter{03BB}{$\lambda$}
\DeclareUnicodeCharacter{03BD}{$\nu$}
\DeclareUnicodeCharacter{03C1}{$\rho$}
\DeclareUnicodeCharacter{03C4}{$\tau$}
\DeclareUnicodeCharacter{03C6}{$\phi$}


%% new maketitle
\makeatletter
\renewcommand*{\maketitle}{%
\begin{titlepage}
 \centering
 \resizebox{\textwidth}{!}{\includegraphics{modulo-head-banner.png}}\\
 \vspace*{\fill}
 {\Huge \fontfamily{Montserrat-TOsF}\selectfont \begin{spacing}{1.5} \@title \end{spacing}}
 \vspace*{3cm}
{\raggedleft\Large\@author}\\
 \vspace*{3cm}
%{\raggedleft\Large\the\day.\the\month.\the\year}
{\raggedleft\Large \@date}
 \vspace*{\fill}\\
 \includegraphics[scale=0.5]{by-nc.eu.png}
\end{titlepage}
}
\makeatother

%% display cards
%\newenvironment{card}{\begin{minipage}{0.5\textwidth}}{\end{minipage}}
\makeatletter
\def\@cardenvironment{minipage}
\def\@tabenvironment{framed}
\renewenvironment{sphinxuseclass}[1]{%
\def\myarg{#1}
\def\cardenv{col-sm-6}
\def\tabenv{tabbed-container}
\ifx\myarg\cardenv\begin{minipage}{0.5\textwidth}\fi
%\ifx\myarg\tabenv\begin{framed}\fi
\ifx\myarg\tabenv\centering\begin{minipage}{0.5\textwidth}\begin{framed}\fi
}%
{%
\ifx\@currenvir\@cardenvironment\end{minipage}\fi
\ifx\@currenvir\@tabenvironment\end{framed}\end{minipage}\fi
%\vspace{0.5cm}
}
\makeatother

%\DeclareStringOption[1pt]{\spx@opt@noteMichaborder}
%\DeclareStringOption[0.5pt]{noteMichaborder}
%\sphinxDeclareSphinxColorOption{noteMichaBorderColor}{{rgb}{0,0,0}}




%\newenvironment{sphinxnoteMicha}[1]{\begin{tcolorbox}[colframe=red]\sphinxstrong{#1}}{\end{tcolorbox}}
%\renewenvironment{sphinxnote}[1]{\begin{sphinxheavybox}\sphinxstrong{#1} }{\end{sphinxheavybox}}

%% personalizing headers and footers (e,g, removing version number)

\setlength{\headheight}{16.6pt}

\definecolor{TitleColor}{rgb}{0, 0, 0}

% customizing Bjornstrup
\ChNameVar{\raggedleft\Large\sffamily\selectfont}
\ChNumVar{\raggedleft\LARGE\sffamily\bfseries\selectfont\vspace{-14pt}}
\ChTitleVar{\vspace{8pt}\centering\LARGE\sffamily\bfseries}
\ChRuleWidth{1.5pt}
\makeatletter

% légende d'images
\usepackage[width=.9\textwidth, font=small, labelfont=bf]{caption}
\DeclareCaptionFont{nohyphen}{\hyphenpenalty=10000 \exhyphenpenalty=10000\relax}
\captionsetup{font+=nohyphen, aboveskip=12pt, belowskip=8pt}

% commencer chaque section sur une nouvelle page impaire
\usepackage{etoolbox}
\preto\section{  
  
\ifnum\value{section}>-1
\cleardoublepage
\vspace{100pt}
\else
\fi
}

\preto\chapter{\cleardoublepage\cleardoublepage}


% éviter les lignes orphelines
\usepackage[%
    all,
    defaultlines=3 % nombre minimum de lignes
]{nowidow}


\renewcommand{\chaptermark}[1]{\markboth{\thechapter.\ #1}{}}
\newcommand{\helv}{%
\fontfamily{phv}\fontsize{9}{11}\selectfont}

\usepackage{datetime}

\makeatletter
   \fancypagestyle{normal}{
    \fancyhf{}
    \renewcommand{\headrulewidth}{0.4pt}
    \renewcommand{\footrulewidth}{0pt}
    \renewcommand{\headruleskip}{3pt}
    \fancyfoot[CO]{{\py@HeaderFamily \mdseries \small \thepage \textcolor{lightgray}{\; $\mid$ \today}}}
    \fancyhead[RO]{{\py@HeaderFamily \mdseries \small \nouppercase{\rightmark}}}
    \if@twoside
     \fancyfoot[CE]{{\py@HeaderFamily \mdseries \small \thepage  \textcolor{lightgray}{\; $\mid$ \today}}}
     \fancyhead[LE]{{\py@HeaderFamily \mdseries \small Modulo \nouppercase{\leftmark}}}
    \fi
    
    
    % define chaptermark with \@chappos when \@chappos is available for Japanese
    \ltx@ifundefined{@chappos}{}
      {\def\chaptermark##1{\markboth{\@chapapp\space\thechapter\space\@chappos\space ##1}{}}}
    }
  % Update the plain style so we get the page number & footer line,
  % but not a chapter or section title.  This is to keep the first
  % page of a chapter `clean.'
   \fancypagestyle{plain}{
    \fancyhf{}
    \fancyfoot[CO]{{\py@HeaderFamily \mdseries \small \thepage \textcolor{lightgray}{\; $\mid$ \today}}}
    \if@twoside\fancyfoot[CE]{{\py@HeaderFamily \mdseries \small \thepage \textcolor{lightgray}{\; $\mid$ \today}}}\fi
    \renewcommand{\headrulewidth}{0pt}
    \renewcommand{\footrulewidth}{0pt}
    \renewcommand{\headruleskip}{3pt}
    }
\makeatother

%%%% in the fly image conversion taken from here: https://tex.stackexchange.com/questions/574562/inclusion-of-graphics-with-on-the-fly-conversion-using-xetex

\usepackage{graphicx}% including pictures
\usepackage{shellesc}% run system commands from TeX engine
\usepackage{etoolbox}% for \csdef, \csletcs
\makeatletter
\def\@firstofthree#1#2#3{#1}%
\def\@secondofthree#1#2#3{#2}%
\newcommand{\DeclareGraphicsAlien}[2]{%
  \edef\Gin@extensions{#1,\Gin@extensions}%
  \DeclareGraphicsRule{#1}{\@gobble#1}{#1}{}%
  \csdef{Gread@\@gobble#1}##1{%
    \edef\SourceFile{\Gin@base\Gin@ext}%
%    \edef\Gin@base{\Gin@base-\@gobble#1-converted-to}%
%    \edef\Gin@base{\Gin@base-\@gobble#1-converted-to}%
    \edef\Gin@ext{#2}%
    \edef\OutputFile{\Gin@base\Gin@ext}%
    \edef\targetfmt{\expandafter\expandafter\expandafter
                    \@firstofthree\csname Gin@rule@\Gin@ext\endcsname\relax}%
    \edef\targetext{\expandafter\expandafter\expandafter
                    \@secondofthree\csname Gin@rule@\Gin@ext\endcsname\relax}%
    %\IfFileExists{\SourceFile}{\ShellEscape{#3}}{}%
    \csletcs{Ginclude@\@gobble#1}{Ginclude@\targetfmt}%
    \csname Gread@\targetfmt\endcsname{\Gin@base\targetext}%
  }%
}%
\makeatother

%%% this needs to be made generalizable

\DeclareGraphicsAlien{.gif}{.pdf}
\DeclareGraphicsAlien{.svg}{.pdf}


